\documentclass{beamer}

\usepackage{graphicx}
\usepackage{amsmath}
\usepackage{amsthm}

% \theoremstyle{definition}
% \newtheorem{definition}{Definition}

% \theoremstyle{plain}
% \newtheorem{theorem}{Theorem}

\theoremstyle{remark}
\newtheorem{remark}{Remark}

\title{Week 5 CUSF State Feedback 2: Linear Quadratic Regulators and }
\author{Yoonjae Hwang}
\date{9 November 2025}

\begin{document}
\frame{\titlepage}

\begin{frame}
\frametitle{Optimal Control}
As we design our controllers, often we are concerned with how to make the \textsl{best} controller for our given system. 
The 'best' control system is dependent on performance metrics we choose, which will be a case by case basis. 
We have to relate a qualitative performance metric to quantitative values. 

\end{frame}

\begin{frame}
\frametitle{Performance vs Effort}
To find a best control parameters, we can define a cost function then solve an optimisation problem to find our parameters that minimises those values. 
To build our cost function, we want it to penalise: 
\begin{itemize}
    \item Low performance
    \item High effort (money, energy, etc)
\end{itemize} 
We want this cost function to be minimised across the entire time in whcih the system operates. 
We also want to be able to weigh these quantities differently, as different scenarios will require different biases. 
\end{frame}

\begin{frame}
\frametitle{The Cost Function}

\begin{equation}
    J = \int_{0}^{\infty} x^{\top} Q x + u^{\top} R u \space dt
\end{equation}

\begin{itemize}
    \item Q and R have to be positive definite matrices. 
    \item The cost function is quadratic - it's convex and has a definite minimum. 
    \item If x and u define a linear system, then the cost function remains quadratic. 
\end{itemize}
\end{frame}

\begin{frame}
\frametitle{Solving the Linear Quadratic Problem}
I won't go into the derivation. 
The end result is: 
\begin{equation}
    u = -K x, \qquad K = R^{-1} B^{\top} S
\end{equation}
where $R, B$ statisfies the \textsl{Algebraic Riccati Equation}
\begin{equation}
    SA + A^{\top} S - S B R^{-1} B^{\top} S + Q = 0
\end{equation}
Remarks: 
\begin{itemize}
    \item Non trivial because it is quadratic in S
    \item single positive-definite solution if and only if the system is stabilizable
    \item You solve the equation numerically. 
\end{itemize}
\end{frame}

\begin{frame}
\frametitle{Integral Action}
State feedback for 2nd order systems are essentially PD controllers - the position variables and their time derivatives are used. 
\newline
To carry out integral control, we must \textsl{augment} the system. We introduce a new term $z = \int_{0}^{t} y(\tau) - r(\tau) d\tau$.
\begin{equation}
    \frac{d}{dt}
    \begin{bmatrix}
        x \\
        z
    \end{bmatrix}
    = 
    \begin{bmatrix}
        Ax + Bu \\
        y - r
    \end{bmatrix}
\end{equation}
The control law is then given by: 
\begin{equation}
    u = - Kx - k_i z + k_f r
\end{equation}
\end{frame}

\begin{frame}
\frametitle{Reachability for Augmented Systems}
The reachability matrix is then given by: 
\begin{equation}
    W_r = 
    \begin{bmatrix}
        B & AB & \cdots & A^n B \\
        0 & CB & \dots & CA^{n-1}B 
    \end{bmatrix}
\end{equation}
If $W_r$ is full rank we say that the system is controllable.
\end{frame}



\end{document}